% problem-000023 2 不稳定含氮化合物反应
\section*{2. 不稳定含氮化合物反应}
某⽩⾊固体A室温时不稳定，⽽在⽔中稳定。A具有弱碱性和较显著的还原性，分⼦只有1个镜⾯，有顺式和反式异构体。在酸性条件下 A 分解为三种产物 B、C 和 D，其中 B 为⼆元化合物，与 $\mathrm { C O _ { 2 } }$ 是等电⼦体。在中性溶液中，A与亚硝酸作⽤可⽣成化合物E和 $\mathbf { C } _ { \circ }$ 。纯净的E是⼀种⽆⾊晶体，为⼆元酸 $( \mathrm { p } K _ { a 1 } = 6 . 9 , \mathrm { p } K _ { a 2 } =$ 11.6)。E 的阴离⼦有 $C _ { 2 }$ 轴和与此轴垂直的镜⾯，其钠盐 F 可通过⾦属钠与硝酸铵1:1 反应⽽成，也可在⼄醇介质中由A与RONO和⼄醇钠反应⽽成。

\subsection*{2-1}
写出以上两种⽅法合成F的反应⽅程式。

\subsection*{2-2}
画出A的顺式和反式异构体的⽴体结构。

\subsection*{2-3}
⼲燥的E晶体极易爆炸，⽔溶液中较稳定，但仍会逐渐分解，写出其分解反应⽅程式。

\subsection*{2-4}
E的异构体G是弱酸， ${ \tt p } K _ { a } = 6 . 6$ ，画出G的路易斯结构。

\subsection*{2-5}
⽤液态 $\mathrm { N _ { 2 } O _ { 4 } }$ 与F反应⽣成H，此法得到的H为�构型，其中Na的百分含量为37.69%。存在⼀个O‒O键。推断H的化学式并画出其阴离⼦的结构（孤对电⼦不需标出）。

\subsection*{2-6}
H的�构型同分异构体（没有O−O键）可由A与 $\mathrm { B u O N O _ { 2 } }$ 和甲醇钠在甲醇介质中反应⽽成，写出�-H的制备反应⽅程式。
